\section{Conclusion}

The compound resilience hole concept reframes the I2.0 climate adaptation investment problem. Prior work \citep{ROUTE_D1} identified D1 climate exposure as a standalone investment driver; prior work \citep{ROUTE_B1} identified B1 structural isolation as a standalone driver. This paper demonstrates that the two dimensions are not independent: a corridor with high D1 and low B1 is weathered but reroutable; a corridor with high B1 and low D1 is isolated but reliable. The compound exposure condition --- high on both --- creates a qualitatively different risk category in which no single mitigation strategy is sufficient.

\subsection*{Key Findings}

\textbf{Compound exposure is concentrated in mountain passes and coastal lowlands.} Of the 11 corridors satisfying the B1 $>$ 7 AND D1 $>$ 6 threshold, seven involve either mountain pass crossings (Donner, Snoqualmie, Siskiyou, Homestake) or Gulf and coastal flood exposure (Gulf Coast I-10 Louisiana and Texas, Southeast I-95 South Carolina). These are not coincidental geography: mountain passes concentrate freight on a single alignment by topographic necessity, and coastal lowlands concentrate freight on elevated causeways with no viable alternate due to bay and estuary geography. Both create the structural conditions for B1 elevation.

\textbf{The Donner Pass freight tunnel is the highest-NPV single investment in the I2.0 program.} At \$4 billion construction cost and \$1.6 billion in annual avoided disruption, the simple payback period is 2.5 years. The 30-year NPV at a 7\% discount rate exceeds \$15 billion. No other single project in the I2.0 program --- not managed freight lanes on I-75 Atlanta, not the I-95 Northeast capacity expansion, not the Gulf Coast hardening program --- matches this return per dollar of investment. The tunnel's NPV superiority is a direct consequence of compound exposure: it simultaneously fixes the D1 component (bypassing the weather-closure segment entirely) and the B1 component (creating a second interstate-grade path through the Sierra Nevada, reducing isolation from B1 = 8.3 to B1 $\approx$ 4.0).

\textbf{Gulf Coast I-10 hardening is the second priority.} The D1 trajectory for Gulf Coast I-10 Louisiana is the steepest in the national corpus under 2050 climate projections \citep{NOAA_SLR2022}: a D1 score rising from 8.4 (2024) to a projected 9.1 (2050) as sea level rise converts Category 1 and 2 storm surge events from nuisance flooding to corridor-closing inundation. At \$2.1 billion for the Louisiana hardening program, the investment has a 30-year NPV exceeding \$8 billion under mid-range climate scenarios. The B1 component of Gulf Coast compound exposure is addressed through designated freight alternates on I-20 and I-12, not through new construction.

\textbf{The compound exposure matrix is the correct tool for prioritizing climate adaptation investment.} The I2.0 climate adaptation program has finite resources. Investments that address D1 alone (corridor hardening) reduce closure frequency but do not eliminate the single-path dependency when closures occur. Investments that address B1 alone (alternate route development) reduce isolation but do not reduce the climate events that trigger the isolation risk. The compound exposure matrix identifies the corridors where both interventions are required, and where the dual-benefit structure of the right intervention (Donner tunnel) provides exceptional return on investment.

\subsection*{Policy Implications}

The finding has direct implications for federal climate resilience funding under IIJA and successor programs. Current FHWA climate resilience programs evaluate investments on D1-equivalent metrics: flood vulnerability, storm surge probability, wildfire exposure. They do not account for network topology (B1) --- the question of whether an at-risk corridor is the only path for a significant share of national freight, or one of several adequate paths. This paper demonstrates that the B1 dimension is not a refinement but an essential input: it determines whether a D1 investment delivers national-scale benefits (Donner, Gulf Coast Louisiana) or corridor-scale benefits (I-95 South Carolina, which has parallel alternates on US-17 and US-1).

The ROUTE compound exposure matrix provides the analytical tool to integrate network topology into climate resilience investment decisions. Future research \citep{ROUTE_C2} will extend this framework to the max-flow consequences of compound exposure: when Donner closes, how does the national freight max-flow change? When Gulf Coast I-10 closes, which O-D pairs lose their primary path and what is the freight cost of the resulting network reconfiguration?

\subsection*{Future Work}

Three extensions are warranted. First, the D1 scores in this analysis use 2024 baseline climate data; 2035 and 2050 projections under IPCC RCP 4.5 and 8.5 scenarios should be incorporated as climate projections are updated. Second, the B1 threshold of 7 is calibrated to the v1.1 corpus; as the ROUTE corpus expands to include T2 and T3 corridors, the compound exposure condition should be recalibrated to the full network distribution. Third, the economic cost model in this paper uses ATRI truck-hour costs \citep{ATRI_costs2024}; a commodity-class-weighted cost model would better capture the asymmetric disruption costs for time-sensitive freight categories (pharmaceuticals, perishables, automotive JIT).
